\section[Cote 68 : les théories de l'orientation et leurs applications de Stokes]{Cote 68 : les théories de l'orientation et leurs applications de Stokes\protect\verifie{Folder68}}\label{sec:68}

La cote 68 réunit des notes sur les jeux positionnels, les réarrangements, l'exponentielle et, aux pages~49 à~55, l'orientation des espaces vectoriels réels et les déterminants. La page~49 (\lot{68}{3}{49}) énonce deux lemmes sur les ensembles à deux éléments, puis les données d'une \q{théorie de l'orientation} sur le groupoïde des espaces vectoriels réels de dimension finie : un foncteur $\Omega$ vers les ensembles à deux éléments, une identification $\Omega(0) \simeq \{\pm 1\}$, et pour tout hyperplan $W \subset V$ et tout demi-espace $V'$ bordé par $W$ une application de Stokes $\mathrm{St}_{V'} : \Omega(V) \to \Omega(W)$, naturelle en $(V, V')$, bijective, qui change par l'involution de $\Omega(W)$ quand on passe au demi-espace opposé. Le théorème de la page dit : a)~tout morphisme de théories est un isomorphisme ; b)~entre deux théories il existe un et un seul isomorphisme ; c)~il existe une théorie. La page démontre a) et l'unicité dans b) par récurrence sur la dimension. La page~51 (\lot{68}{3}{51}) ramène une théorie en dimension $n$ à un caractère $\varepsilon_n$ de $\mathrm{GL}(n, \mathbf R)$ et à une bijection de Stokes $\omega_n \to \omega_{n-1}$ pour le demi-espace $\mathbf R^{n-1} \times \mathbf R_+$, écrit la compatibilité au stabilisateur $\varepsilon_n(u_n) = \mathrm{sg}(\lambda)\,\varepsilon_{n-1}(u_{n-1})$ et conclut \q{de proche en proche} $\varepsilon_n = \mathrm{sg} \circ \det$. On démontre ici a), b) et la valeur $\mathrm{sg} \circ \det$ de l'action, c'est-à-dire tout l'énoncé de la piste ; l'existence c) n'en fait pas partie.

\noindent\emph{Conventions.} Les espaces sont les espaces vectoriels réels de dimension finie (dans l'univers \texttt{Type}) et leurs isomorphismes linéaires. Un demi-espace fermé bordé par $W$ est un ensemble $H = \{f \geq 0\}$, où $f$ est une forme linéaire non nulle de noyau $W$ ; le demi-espace opposé est $-H = \{f \leq 0\}$. Une théorie de l'orientation est la donnée d'un ensemble $\Omega(V)$ à deux éléments pour chaque $V$, d'une action fonctorielle des isomorphismes $e : V \to V_2$, d'une bijection $\Omega(\mathbf R^0) \simeq \{\pm 1\}$ et, pour chaque couple $(W, H)$, d'une application $\mathrm{St}_H : \Omega(V) \to \Omega(W)$, soumise aux conditions suivantes. Naturalité : si $e : V \to V_2$ envoie $W$ sur $W_2$ et $H$ sur $H_2$, alors $\mathrm{St}_{H_2} \circ e = e|_W \circ \mathrm{St}_H$. Bijectivité de $\mathrm{St}_H$. Opposé : $\mathrm{St}_{-H}(\omega) \neq \mathrm{St}_H(\omega)$ pour tout $\omega$, ce qui, dans un ensemble à deux éléments, revient à $\mathrm{St}_{-H} = \tau \circ \mathrm{St}_H$, où $\tau$ est l'involution sans point fixe. Un morphisme de théories est une transformation naturelle $\varphi_V : \Omega(V) \to \Omega'(V)$ compatible à l'identification de $\Omega(0)$ et aux applications de Stokes. On note $\mathbf R^n_+ = \mathbf R^{n-1} \times \mathbf R_+$ le demi-espace standard de $\mathbf R^n$, bordé par $\mathbf R^{n-1} \times 0$.

\begin{lemme}\label{lem:68-deux}
Soient $P$, $Q$ des ensembles à deux éléments. (1) Deux éléments distincts d'un même troisième sont égaux. (2) Deux applications injectives $P \to Q$ qui coïncident en un point sont égales. (3) Une bijection de $P$ fixe tous les points ou n'en fixe aucun ; la composée de deux bijections fixe tous les points si et seulement si les deux fixent tous les points ou aucune n'en fixe ; le carré d'une bijection fixe tous les points.
\end{lemme}

\begin{proof}
(1) est la définition du cardinal~$2$. Pour (2), si $y \neq x_0$, $f(y)$ et $g(y)$ sont distincts de $f(x_0) = g(x_0)$, donc égaux par~(1). Pour (3), si $\sigma(x) \neq x$ et $\sigma(y) = y$, alors $x \neq y$ et $\sigma(x) \neq \sigma(y) = y$, donc $\sigma(x) = x$ par~(1), contradiction. Si ni $\sigma$ ni $\rho$ ne fixent de point, $\sigma(\rho x)$ et $x$ sont distincts de $\rho x$, donc égaux ; les autres cas sont immédiats.
\end{proof}

\begin{enonce}[Page 51, l'action de $\mathrm{GL}(n, \mathbf R)$\piste{68-orientation-stokes-axioms}]
Soit $(\Omega, \mathrm{St})$ une théorie de l'orientation et $g \in \mathrm{GL}(n, \mathbf R)$. Pour tout $\omega \in \Omega(\mathbf R^n)$, $g\omega = \omega$ si $\det g > 0$ et $g\omega = \tau\omega$ si $\det g < 0$ : l'action de $\mathrm{GL}(n, \mathbf R)$ sur $\Omega(\mathbf R^n)$ est $\mathrm{sg} \circ \det$. Il en est de même pour un automorphisme de tout espace $V$.
\end{enonce}

\begin{proof}
Disons qu'une matrice inversible $M$ est \emph{bonne} si l'automorphisme de matrice $M$ fixe $\Omega(\mathbf R^n)$ si et seulement si $\det M > 0$. Par le lemme~\ref{lem:68-deux}~(3), un produit de deux matrices bonnes est bon, car $\det(AB) > 0$ si et seulement si $\det A$ et $\det B$ sont de même signe, et un carré $A^2$ est bon, car il fixe $\Omega(\mathbf R^n)$ et $\det A^2 > 0$. Une transvection $1 + c\,E_{ij}$ est le carré de $1 + \tfrac c2 E_{ij}$ ; elle est bonne. Une matrice diagonale $\mathrm{diag}(1, \ldots, d, \ldots, 1)$ avec $d > 0$ est le carré de $\mathrm{diag}(1, \ldots, \sqrt d, \ldots, 1)$ ; elle est bonne.

Reste $r_i = \mathrm{diag}(1, \ldots, -1, \ldots, 1)$, avec $-1$ en place $i$, dont le déterminant est $-1$ : il faut voir qu'elle ne fixe pas $\Omega(\mathbf R^n)$. Pour $i = n$, $r_n$ fixe l'hyperplan $\mathbf R^{n-1} \times 0$ point par point et envoie $\mathbf R^n_+$ sur le demi-espace opposé. La naturalité donne $\mathrm{St}_{-\mathbf R^n_+}(r_n\omega) = \mathrm{St}_{\mathbf R^n_+}(\omega)$ ; si $r_n\omega = \omega$, cela contredit l'axiome de l'opposé. Pour $i$ quelconque, $r_i$ est conjuguée de $r_n$ par la matrice de la transposition de $i$ et $n$, et un conjugué fixe $\Omega$ si et seulement si l'élément conjugué le fait. Une matrice diagonale inversible est un produit de matrices $\mathrm{diag}(1, \ldots, d_i, \ldots, 1)$, et pour $d_i < 0$ on écrit $\mathrm{diag}(\ldots, d_i, \ldots) = \mathrm{diag}(\ldots, -d_i, \ldots)\,r_i$ : toute matrice diagonale inversible est bonne.

Toute matrice inversible est produit de matrices diagonales inversibles et de transvections (par un principe de récurrence de mathlib) ; elle est donc bonne. Si $g$ ne fixe pas $\Omega(\mathbf R^n)$, il n'en fixe aucun point (lemme~\ref{lem:68-deux}~(3)), et $g\omega = \tau\omega$. Pour un espace $V$ et un isomorphisme $c : V \to \mathbf R^n$, $g$ fixe $\Omega(V)$ si et seulement si $c g c^{-1}$ fixe $\Omega(\mathbf R^n)$, et $\det(c g c^{-1}) = \det g$.
\end{proof}

\begin{encadre}
\begin{theoreme}[Grothendieck, page 49, théorème, a) et unicité dans b)\piste{68-orientation-stokes-axioms}]
Soient $(\Omega, \mathrm{St})$ et $(\Omega', \mathrm{St}')$ deux théories de l'orientation. (1) Toutes les composantes d'un morphisme $\varphi$ sont bijectives, et $\varphi$ a un inverse qui est un morphisme de théories. (2) Deux morphismes $\varphi, \psi : \Omega \to \Omega'$ sont égaux.
\end{theoreme}
\end{encadre}

\begin{proof}
Par récurrence sur $n = \dim V$, on montre que $\varphi_V$ est injective et que $\varphi_V = \psi_V$. Si $n = 0$, on choisit un isomorphisme $c : V \to \mathbf R^0$. Par naturalité et compatibilité à $\Omega(0) \simeq \{\pm 1\}$, $\varphi_V$ et $\psi_V$ s'écrivent toutes deux $c^{-1} \circ \iota'^{-1} \circ \iota \circ c$, où $\iota$ et $\iota'$ sont les deux identifications ; elles sont égales et injectives. Si $n = m + 1$, on choisit une forme linéaire non nulle $f$ ; son noyau $W$ est de dimension $m$, et $H = \{f \geq 0\}$ est un demi-espace bordé par $W$. La compatibilité aux applications de Stokes donne $\mathrm{St}'_H \circ \varphi_V = \varphi_W \circ \mathrm{St}_H$, et de même pour $\psi$. Comme $\varphi_W = \psi_W$ est injective par hypothèse de récurrence et que $\mathrm{St}_H$ et $\mathrm{St}'_H$ sont bijectives, $\varphi_V$ est injective et $\mathrm{St}'_H \circ \varphi_V = \mathrm{St}'_H \circ \psi_V$, d'où $\varphi_V = \psi_V$. Une injection entre deux ensembles à deux éléments est bijective. Pour l'inverse, on prend le morphisme $\Omega' \to \Omega$ de l'énoncé suivant ; ses composées avec $\varphi$ sont des morphismes d'une théorie dans elle-même, donc égales à l'identité par~(2).
\end{proof}

\begin{enonce}[Page 49, théorème, existence dans b)\piste{68-orientation-stokes-axioms}]
Entre deux théories de l'orientation $(\Omega, \mathrm{St})$ et $(\Omega', \mathrm{St}')$, il existe un morphisme ; c'est donc l'unique isomorphisme entre elles.
\end{enonce}

\begin{proof}
On suit la réduction de la page~51. Notons $\omega_n = \Omega(\mathbf R^n)$ et $s_n : \omega_n \to \omega_{n-1}$ la bijection composée de $\mathrm{St}_{\mathbf R^n_+}$ et de l'identification $\mathbf R^{n-1} \times 0 \simeq \mathbf R^{n-1}$. On définit des éléments de référence $r_0 \in \omega_0$, image réciproque de $1$, et $r_n = s_n^{-1}(r_{n-1})$ ; de même $r'_n \in \omega'_n$. Soit $\psi_n : \omega_n \to \omega'_n$ la bijection qui envoie $r_n$ sur $r'_n$ et l'autre élément sur l'autre. Elle commute à l'action de $\mathrm{GL}(n, \mathbf R)$ : par l'énoncé de la page~51, $g$ agit trivialement des deux côtés si $\det g > 0$, et par $\tau$ des deux côtés si $\det g < 0$, et une bijection commute à $\tau$ (lemme~\ref{lem:68-deux}~(1)). Par le lemme~\ref{lem:68-deux}~(2), elle vérifie $s'_n \circ \psi_n = \psi_{n-1} \circ s_n$, puisque les deux membres sont injectifs et envoient $r_n$ sur $r'_{n-1}$, et elle est compatible à l'identification de $\Omega(0)$.

Pour un espace $V$ de dimension $n$ et un isomorphisme $c : V \to \mathbf R^n$, on pose $\varphi_V = c^{-1} \circ \psi_n \circ c$. Cela ne dépend pas de $c$ : si $c_2 = g c$ avec $g \in \mathrm{GL}(n, \mathbf R)$, $\psi_n$ commute à $g$. La naturalité de $\varphi$ s'en déduit, en calculant $\varphi_{V_2}$ avec la carte $c \circ e^{-1}$. Pour la compatibilité à $\mathrm{St}_H$, $H = \{f \geq 0\}$ bordé par $W$ dans $V$ de dimension $m + 1$, on construit un isomorphisme $c : V \to \mathbf R^{m+1}$ dont la dernière coordonnée est $f$ : avec $v_0$ tel que $f(v_0) = 1$ et un isomorphisme $j : W \to \mathbf R^m$, l'application $u \mapsto j^{-1}(u_1, \ldots, u_m) + u_{m+1} v_0$ est injective entre espaces de même dimension, et $c$ est son inverse. Alors $c$ envoie $W$ sur $\mathbf R^m \times 0$ et $H$ sur $\mathbf R^{m+1}_+$. La naturalité des applications de Stokes pour $c$ et pour $c^{-1}$ ramène $\mathrm{St}'_H \circ \varphi_V = \varphi_W \circ \mathrm{St}_H$ à la relation $s'_{m+1} \circ \psi_{m+1} = \psi_m \circ s_{m+1}$, $\varphi_W$ étant calculé avec la carte $j' \circ c|_W$, où $j'$ identifie $\mathbf R^m \times 0$ à $\mathbf R^m$.
\end{proof}

\begin{remarque}
La vérification suit la lecture, qui pose la bijectivité des applications de Stokes en axiome : l'argument de cardinal de la page~49 ne la donne pas, et elle sert à chaque pas de la récurrence. La partie a) et l'unicité dans b) n'utilisent ni l'axiome de l'opposé ni la naturalité des applications de Stokes : elles n'utilisent que leur bijectivité, la compatibilité des morphismes et l'identification de $\Omega(0)$. L'axiome de l'opposé ne sert qu'à l'existence, par l'action de $\mathrm{diag}(1, \ldots, 1, -1)$ ; sans lui, la théorie où $\Omega(V) = \{\pm 1\}$, avec action triviale et applications de Stokes identiques, satisferait aux autres axiomes sans être isomorphe aux orientations usuelles (observation non formalisée). Le passage \q{de proche en proche} de la page~51 repose, d'après la note de la lecture, sur la factorisation par le déterminant de tout caractère de $\mathrm{GL}(n, \mathbf R)$ à valeurs dans $\{\pm 1\}$ ; la démonstration tire cette factorisation de la décomposition des matrices inversibles en transvections et matrices diagonales, les transvections et les matrices diagonales positives étant des carrés. La forme \q{$\Omega(W, V) \wedge \Omega(V) \simeq \Omega(W)$} des axiomes et le lemme~2 de la page~49 ne sont pas formalisés : on énonce directement l'axiome de l'opposé.
\end{remarque}

Ne sont pas démontrés ici : l'existence c) d'une théorie de l'orientation, c'est-à-dire la construction des applications de Stokes sur les orientations usuelles et la vérification des axiomes ; la reformulation par le produit contracté $D(W, V) \wedge \Omega(V) \simeq \Omega(W)$ et le lemme~2 de la page~49 ; la comparaison de la normalisation de la page (normale rentrante en dernier) avec celle des géomètres, au signe $(-1)^{\dim V}$ ; le produit, la règle de Koszul et la compatibilité de Stokes et des produits de la page~53 ; les déterminants de la page~55.
